\documentclass[11pt]{article}
\usepackage[a4paper,margin=1.9cm]{geometry}
\usepackage[T1]{fontenc}
\usepackage{textcomp}
\usepackage[spanish]{babel}
\usepackage{amsmath,amssymb}
\usepackage{enumitem}
\usepackage{tikz}
\usetikzlibrary{calc,arrows.meta,patterns,decorations.pathmorphing}
\usepackage{xcolor}
\usepackage{multicol}
\usepackage{parskip}
\usepackage{lmodern}
\renewcommand{\familydefault}{\sfdefault}

\definecolor{unalblue}{RGB}{31,73,124}
\definecolor{unalblue2}{RGB}{70,120,180}

\newlist{opciones}{enumerate}{1}
\setlist[opciones]{label=\textbf{(\alph*)},itemsep=1pt,topsep=2pt,leftmargin=1.8em,font=\normalfont}
\setlist[enumerate,1]{label=\textbf{\arabic*.},itemsep=6pt,topsep=4pt,leftmargin=1.6em,font=\normalfont}
\setlist[enumerate,2]{label=\textbf{\alph*)},itemsep=4pt,topsep=3pt,leftmargin=1.6em,font=\normalfont}
\setlist[enumerate,3]{label=\textbf{\roman*)},itemsep=2pt,topsep=2pt,leftmargin=1.4em,font=\normalfont}

\newcommand{\vect}[2]{\begin{pmatrix}#1\\#2\end{pmatrix}}
\newcommand{\vectiii}[3]{\begin{pmatrix}#1\\#2\\#3\end{pmatrix}}

\newcommand{\examheader}{%
\noindent\begin{tikzpicture}
\node[fill=unalblue, text width=\dimexpr\linewidth-16pt\relax, inner sep=8pt, rounded corners=3pt, align=left, text=white] (hdr) {%
  \begin{minipage}[c]{0.66\linewidth}
    {\Large\bfseries \examsubject}\\[2pt]
    {\small \examtype\ \textbullet\ \textcolor{white!85!unalblue}{\examsubtitle}}
  \end{minipage}\hfill
  \begin{minipage}[c]{0.28\linewidth}
    \raggedleft{\scriptsize Escuela de Matem\'aticas\\[-1pt] Universidad Nacional\\[-1pt] de Colombia, Sede Medell\'in}
  \end{minipage}%
};
\end{tikzpicture}\\[7pt]
}
\newcommand{\examfooter}{%
  \vfill
  \rule{\linewidth}{0.4pt}\\[1pt]
  {\scriptsize\textit{Transcripci\'on digital --- MathUNAL}\hfill\textcolor{unalblue}{\textbf{\thepage}}}
}
\pagestyle{empty}

\newcommand{\examsubject}{ECUACIONES DIFERENCIALES}
\newcommand{\examtype}{Tercer Parcial}
\newcommand{\examsubtitle}{Semestre 02-2019, 22 de Agosto de 2019}

\begin{document}

\examheader

\vspace{4pt}
\noindent
\begin{minipage}[t]{0.55\linewidth}
{\bfseries\large Tercer Parcial de Ecuaciones Diferenciales}\\[2pt]
Puntaje. S\'olo para uso Oficial\\[4pt]
\begin{tabular}{|c|c|c|c|c|c|c|}
\hline
1 (25) & 2 (10) & 3 (25) & 4 (10) & 5 (30) & Total & Nota\\
\hline
 & & & & & & \\
\hline
\end{tabular}
\end{minipage}%
\hfill
\begin{minipage}[t]{0.4\linewidth}
Nombre: \dotfill\\[6pt]
C\'edula: \dotfill \hfill Grupo: \dotfill\\[10pt]
{\large Examen N\'umero: \dotfill}\\[10pt]
Firma: \dotfill
\end{minipage}

\vspace{6pt}
\noindent\textbf{Instrucciones:} La duraci\'on del examen es de 1 hora y 45 minutos. El examen consta de cuatro preguntas en dos hojas impresas por ambos lados, antes de empezar verifique que su examen est\'e \textbf{completo} y que sus \textbf{datos} sean \textbf{correctos}. No \textbf{desengrape} su examen ni \textbf{a\~nada} otras grapas o su examen ser\'a \textbf{anulado}. En las preguntas con procedimiento \textbf{justifique} sus respuestas en los \textbf{espacios asignados}. No est\'a permitido hacer preguntas, usar calculadoras, sacar hojas en blanco ni ning\'un tipo de apuntes durante el examen. Por favor verifique que su celular est\'e \textbf{apagado}.

\examfooter
\newpage

\examheader

\begin{enumerate}
\item \textbf{(25\%)}
\begin{enumerate}[label=\alph*)]
\item \textbf{(15\%)} Encuentre la soluci\'on general del sistema de ecuaciones lineales de primer orden
\[
x' = x + y, \qquad y' = 4x + y.
\]
\item \textbf{(5\%)} ¿Es verdad que
\[
\lim_{t\to\infty}\begin{pmatrix} x(t) \\ y(t) \end{pmatrix} = \begin{pmatrix} 0 \\ 0 \end{pmatrix}
\]
para toda soluci\'on del sistema anterior? Justifique su respuesta.
\item \textbf{(5\%)} Si el sistema describe una part\'icula que en el tiempo $t=0$ se encuentra en la posici\'on $(0,2)$, encuentre la posici\'on $(x(t),y(t))$ de la part\'icula para todo tiempo $t$.
\end{enumerate}

\vspace{8pt}
\item \textbf{(10\%)} Dados $a,b \in \mathbb{R}$ determinar
\[
\mathcal{L}\left(\int_0^t e^{a\tau}\,d\tau \;*\; \int_0^t e^{b\tau}\,d\tau\right).
\]
\end{enumerate}

\examfooter
\newpage

\examheader

\begin{enumerate}
\setcounter{enumi}{2}
\item \textbf{(25\%)} Determinar
\[
\mathcal{L}^{-1}\left(\frac{(s-2)e^{-s}}{s^2+4s+13}\right).
\]

\vspace{8pt}
\item \textbf{(10\%)} Resuelva la ecuaci\'on integral
\[
y(t) = \int_0^t y(t) e^{t-\tau}\,d\tau + 2\delta(t).
\]
\end{enumerate}

\examfooter
\newpage

\examheader

\begin{enumerate}
\setcounter{enumi}{4}
\item \textbf{(30\%)} En los literales a), b), c) y d) complete seg\'un corresponda. S\'olo se calificar\'a respuesta, no es necesario que escriba el procedimiento. \textbf{Simplifique su respuesta.}
\begin{enumerate}[label=\alph*)]
\item \textbf{(5\%)} El vector $X = k\begin{pmatrix} 4 \\ 5 \end{pmatrix}$, donde $k$ es una constante, es una soluci\'on del sistema
\[
\begin{pmatrix} x \\ y \end{pmatrix}' = \begin{pmatrix} 1 & 4 \\ 2 & -1 \end{pmatrix}\begin{pmatrix} x \\ y \end{pmatrix} - \begin{pmatrix} 8 \\ 1 \end{pmatrix}
\]
para $k = $ \dotfill .
\item \textbf{(5\%)} Si $\mathcal{L}(f) = F$ y $\mathcal{L}(g) = G$, entonces $\mathcal{L}^{-1}(FG) = $ \dotfill .
\item \textbf{(10\%)}
\[
\mathcal{L}\left(\int_0^t e^{a\tau} f(\tau)\,d\tau\right) = \text{\dotfill},
\]
mientras que
\[
\mathcal{L}\left(e^{a\tau}\int_0^t f(\tau)\,d\tau\right) = \text{\dotfill}.
\]
\item \textbf{(10\%)} Determine la soluci\'on general del sistema de ecuaciones
\[
\begin{pmatrix} x \\ y \end{pmatrix}' = \begin{pmatrix} 4 & -2 \\ 8 & -4 \end{pmatrix}\begin{pmatrix} x \\ y \end{pmatrix} + \begin{pmatrix} t^{-3} \\ -t^{-2} \end{pmatrix},
\]
asumiendo que la matriz del sistema tiene un \'unico valor propio $r=0$, y que todo vector propio es m\'ultiplo del vector $(1,2)$.
\[
\text{\dotfill}
\]
\end{enumerate}
\end{enumerate}

\vspace{4pt}
{\scriptsize\textit{Nota de transcripci\'on: el examen original numera este \'ultimo punto como ``4.'' por error de imprenta (repite el n\'umero del punto anterior) — se renumer\'o como punto 5 para que coincida con la tabla de puntaje (1, 2, 3, 4, 5) impresa en la portada.}}

\examfooter

\end{document}
